\documentclass[12pt,a4paper]{article}
\usepackage[utf8]{inputenc}
\usepackage[T1]{fontenc}
\usepackage[french]{babel}
\usepackage{geometry}
\usepackage{fancyhdr}
\usepackage{amsmath}
\usepackage{listings}
\usepackage{xcolor}
\usepackage{hyperref}
\usepackage{graphicx}
\usepackage{enumitem}

\geometry{margin=2.5cm}
\setlength{\headheight}{14.5pt}

\pagestyle{fancy}
\fancyhf{}
\fancyhead[L]{R\'esum\'e Administration Unix}
\fancyhead[R]{\thepage}
\fancyfoot[C]{}

\lstset{
    language=bash,
    basicstyle=\ttfamily\small,
    keywordstyle=\color{blue},
    commentstyle=\color{gray},
    stringstyle=\color{red},
    numbers=left,
    numberstyle=\tiny\color{gray},
    breaklines=true,
    frame=single,
    showstringspaces=false
}

\title{\textbf{R\'esum\'e Complet - Administration Unix}}
\author{}
\date{}

\begin{document}

\maketitle
\thispagestyle{fancy}

\section*{Note importante}
Les parties suivantes sont exclues de ce r\'esum\'e :
\begin{itemize}
    \item L'int\'egration dans les r\'eseaux (TP5)
    \item LVM (TP8)
    \item RAID (TP9)
    \item Sauvegarde et restauration (TP10)
    \item Quotas
\end{itemize}

\newpage
\tableofcontents
\newpage

\section{Introduction et G\'en\'eralit\'es}

\subsection{Histoire d'Unix}
\begin{itemize}
    \item \textbf{1969} : Cr\'eation d'UNIX par Ken Thompson aux laboratoires Bell
    \item \textbf{1973} : R\'e\'ecriture compl\`ete en langage C par Denis Ritchie
    \item \textbf{1974} : AT\&T propose les premi\`eres licences aux universit\'es
    \item \textbf{Ann\'ees 80} : Clonage d'UNIX par d'autres constructeurs (ULTRIX, BSD, AIX)
    \item \textbf{1983} : Lancement du projet GNU par Richard Stallman
    \item \textbf{1991} : Cr\'eation du noyau Linux par Linus Torvalds
\end{itemize}

\subsection{Linux : D\'efinition}
Linux est le noyau d'un OS de type Unix, initi\'e par Linus Torvalds en 1991. Il d\'esigne un OS comparable \`a Windows ou Mac OS compos\'e de nombreux logiciels libres.

\subsection{Caract\'eristiques de Linux}
\begin{itemize}
    \item Code source disponible (licence GPL)
    \item Syst\`eme multit\^ache et multi-utilisateur
    \item Multi-plateforme (Intel x86, Sun Sparc, etc.)
    \item Gestion du multiprocesseur (SMP)
    \item Compatible POSIX
    \item Gestion des consoles virtuelles
    \item Impl\'ementation compl\`ete de la pile TCP/IP
    \item Interface graphique X-Window
\end{itemize}

\subsection{Architecture d'un syst\`eme Linux}
Trois couches :
\begin{enumerate}
    \item \textbf{Couche physique} : p\'eriph\'eriques + BIOS
    \item \textbf{Couche syst\`eme} : kernel et processus
    \item \textbf{Couche interface} : shell + X-Window
\end{enumerate}

\subsection{Le Shell}
Le Shell est l'interpr\'eteur de commande, interface entre l'utilisateur et l'OS. Types de shell disponibles : bash (sh), csh, ksh, zsh, etc.

\section{Standard de Hi\'erarchie des Syst\`emes de Fichiers (FHS)}

\subsection{D\'efinition}
Le FHS (Filesystem Hierarchy Standard) d\'efinit l'arborescence et le contenu des principaux r\'epertoires des syst\`emes de fichiers Linux/Unix.

\subsection{R\'epertoires principaux}

\subsubsection{/bin, /sbin, /lib}
\begin{itemize}
    \item \textbf{/bin} : ex\'ecutables pour tous les utilisateurs (ls, cp, mv, vi, bash)
    \item \textbf{/sbin} : ex\'ecutables pour administration (shutdown, ifconfig, fsck)
    \item \textbf{/lib} : biblioth\`eques partag\'ees (shared libraries)
\end{itemize}

\subsubsection{/dev}
Contient des fichiers sp\'eciaux (device files) correspondant aux p\'eriph\'eriques.
\begin{itemize}
    \item Block device files (disques)
    \item Character device files (terminaux)
    \item Commande : \texttt{mknod} pour cr\'eer des fichiers sp\'eciaux
\end{itemize}

\subsubsection{/etc}
Contient les fichiers et scripts de configuration des diff\'erents services :
\begin{itemize}
    \item /etc/X11 : configuration X-window
    \item /etc/rc.d : scripts de d\'emarrage
    \item /etc/cron : t\^aches p\'eriodiques
    \item /etc/skel : fichiers \`a recopier pour nouveaux utilisateurs
    \item /etc/sysconfig : configuration des p\'eriph\'eriques
\end{itemize}

\subsubsection{/usr}
Unix System Resources - programmes et utilitaires non indispensables :
\begin{itemize}
    \item /usr/bin, /usr/sbin, /usr/lib : \'equivalents de /bin, /sbin, /lib
    \item /usr/etc : fichiers de configuration (rarement utilis\'e)
    \item /usr/include : fichiers .h pour le compilateur C
    \item /usr/local : arborescence propre \`a la machine
    \item /usr/share : fichiers ind\'ependants de l'architecture
    \item /usr/src : sources
\end{itemize}

\subsubsection{/var}
Fichiers dont la taille peut cro\^itre :
\begin{itemize}
    \item /var/lock : fichiers de verrouillage
    \item /var/log : journaux du syst\`eme
    \item /var/spool : files d'attentes (cron, lpd, mail)
    \item /var/run : fichiers contenant les PID des services
\end{itemize}

\subsubsection{/proc}
Pseudo-syst\`eme de fichiers utilis\'e comme interface avec les structures de donn\'ees du noyau. Fichiers virtuels cr\'e\'es dynamiquement en m\'emoire.

\subsubsection{Autres r\'epertoires}
\begin{itemize}
    \item \textbf{/boot} : fichiers utiles pour le chargeur et noyaux Linux
    \item \textbf{/home} : espaces priv\'es des utilisateurs
    \item \textbf{/mnt} : points de montage des partitions externes
    \item \textbf{/media} : points de montage des unit\'es amovibles
    \item \textbf{/tmp} : fichiers temporaires
    \item \textbf{/root} : espace de travail de l'administrateur
    \item \textbf{/lost+found} : fichiers perdus et r\'ecup\'er\'es par fsck
\end{itemize}

\section{D\'emarrage et Arr\^et du Syst\`eme Linux}

\subsection{Processus de d\'emarrage}
\begin{enumerate}
    \item Power on
    \item BIOS/UEFI
    \item BootLoader (GRUB)
    \item Linux Kernel
    \item Processus PID 1 (init/systemd)
    \item System Ready
\end{enumerate}

\subsection{BIOS vs UEFI}
\begin{itemize}
    \item \textbf{BIOS} : Basic Input Output System, situ\'e sur la ROM de la carte m\`ere
    \item \textbf{UEFI} : Unified Extensible Firmware Interface, remplace le BIOS
    \begin{itemize}
        \item Interface graphique \`a fen\^etre
        \item Acc\`es \`a Internet
        \item Support GPT (partitions $>$ 2TB)
        \item Architecture modulaire
    \end{itemize}
\end{itemize}

\subsection{Chargeur de d\'emarrage (BootLoader)}
\begin{itemize}
    \item \textbf{GRUB} : GRand Unified Bootloader (version 2)
    \item Fichiers GRUB2 :
    \begin{itemize}
        \item /etc/default/grub : param\`etres du menu
        \item /etc/grub.d/ : scripts de cr\'eation du menu
        \item /boot/grub2/grub.cfg : fichier de configuration final (g\'en\'er\'e automatiquement)
    \end{itemize}
    \item Commande : \texttt{grub2-mkconfig} pour r\'eg\'en\'erer grub.cfg
\end{itemize}

\subsection{InitRamFs (Initial RAM Disk)}
\begin{itemize}
    \item Syst\`eme de fichiers racine temporaire
    \item Contient les pilotes n\'ecessaires pour acc\'eder au syst\`eme de fichiers racine
    \item Charg\'e en m\'emoire par le bootloader
    \item Permet de monter le "vrai" syst\`eme de fichiers racine
\end{itemize}

\subsection{Gestionnaires de services}

\subsubsection{SysVinit}
\begin{itemize}
    \item Processus /sbin/init (PID = 1)
    \item Niveaux d'ex\'ecution (runlevels) :
    \begin{itemize}
        \item 0 : arr\^et
        \item 1 : mode utilisateur unique (maintenance)
        \item 3 : mode multi-utilisateurs complet (texte)
        \item 5 : mode multi-utilisateurs graphique
        \item 6 : red\'emarrage
    \end{itemize}
    \item Fichier /etc/inittab : configuration
    \item R\'epertoires /etc/rc.d/rcX.d : scripts par niveau
\end{itemize}

\subsubsection{Systemd}
\begin{itemize}
    \item Remplace /sbin/init
    \item Notion d'unit\'e (service, target, mount, socket)
    \item R\'epertoires :
    \begin{itemize}
        \item /usr/lib/systemd/system : unit\'es syst\`eme par d\'efaut
        \item /etc/systemd/system : unit\'es personnalis\'ees
    \end{itemize}
    \item Commandes principales :
    \begin{itemize}
        \item \texttt{systemctl status service} : \'etat d'un service
        \item \texttt{systemctl start/stop/restart service} : contr\^ole
        \item \texttt{systemctl enable/disable service} : activation au boot
        \item \texttt{systemctl mask/unmask service} : bloquer/d\'ebloquer
    \end{itemize}
    \item Cibles (targets) remplacent les runlevels :
    \begin{itemize}
        \item poweroff.target (runlevel 0)
        \item rescue.target (runlevel 1)
        \item multi-user.target (runlevel 3)
        \item graphical.target (runlevel 5)
        \item reboot.target (runlevel 6)
    \end{itemize}
\end{itemize}

\subsection{Arr\^et du syst\`eme}
\begin{itemize}
    \item \texttt{systemctl poweroff} : arr\^et et hors-tension
    \item \texttt{systemctl halt} : arr\^et sans hors-tension
    \item \texttt{systemctl reboot} : red\'emarrage
    \item \texttt{systemctl suspend} : suspension (RAM)
    \item \texttt{systemctl hibernate} : hibernation (disque)
    \item \texttt{shutdown [OPTION] [HEURE] [MESSAGE]} : arr\^et programm\'e
\end{itemize}

\subsection{Gestion des services}
\begin{itemize}
    \item \texttt{systemctl status service} : \'etat d\'etaill\'e
    \item \texttt{systemctl is-active service} : v\'erifier si actif
    \item \texttt{systemctl is-enabled service} : v\'erifier activation au boot
    \item \texttt{systemctl daemon-reload} : recharger configuration
    \item \texttt{journalctl -u service -f} : logs en temps r\'eel
\end{itemize}

\subsection{D\'epannage et journalisation}
\begin{itemize}
    \item \texttt{journalctl} : afficher tous les logs
    \item \texttt{journalctl -f} : suivre les nouveaux logs
    \item \texttt{journalctl -b} : logs du d\'emarrage actuel
    \item \texttt{journalctl -b -1} : logs du d\'emarrage pr\'ec\'edent
    \item \texttt{journalctl -u service} : logs d'un service
    \item \texttt{journalctl -p err} : logs d'erreur et sup\'erieurs
    \item \texttt{journalctl --since "1 hour ago"} : logs depuis 1 heure
\end{itemize}

\section{Gestion des Paquetages}

\subsection{Notion de paquetage}
Un paquet est un fichier archive contenant le logiciel \`a installer et des r\`egles d'installation :
\begin{itemize}
    \item Gestion des d\'ependances
    \item Pr\'e-installation et post-installation
    \item Informations compl\'ementaires (nom, version, description)
\end{itemize}

\subsection{Formats de paquetages}
\begin{itemize}
    \item \textbf{tar} : format ancien, sources et binaires
    \item \textbf{rpm} : RedHat Package Manager (RedHat, SUSE, Fedora)
    \item \textbf{deb} : Debian Package (Debian, Ubuntu, Mint)
\end{itemize}

\subsection{Paquetage RPM}
\begin{itemize}
    \item Outil RPM : Redhat Package Manager
    \item Base de donn\'ees : /var/lib/rpm/
    \item Commandes principales :
    \begin{itemize}
        \item \texttt{rpm -i fichier\_paquetage} : installer
        \item \texttt{rpm -e nom\_du\_paquetage} : d\'esinstaller
        \item \texttt{rpm -U fichier\_paquetage} : mettre \`a jour
        \item \texttt{rpm -qa} : lister tous les packages install\'es
        \item \texttt{rpm -q package} : version install\'ee
        \item \texttt{rpm -qi package} : informations sur un package
        \item \texttt{rpm -ql package} : liste des fichiers install\'es
        \item \texttt{rpm -qf nom\_du\_fichier} : package auquel appartient un fichier
        \item \texttt{rpm -V package} : v\'erification
        \item \texttt{rpm -K package.rpm} : v\'erifier signature
    \end{itemize}
\end{itemize}

\subsection{Gestionnaire de paquets YUM}
Yellowdog Updater Modifier - surcouche de RPM g\'erant t\'el\'echargements et d\'ependances.

Commandes principales :
\begin{itemize}
    \item \texttt{yum list all} : lister tous les paquetages disponibles
    \item \texttt{yum list installed} : lister les paquetages install\'es
    \item \texttt{yum info package} : informations sur un paquetage
    \item \texttt{yum search mot\_cl\'e} : rechercher des paquetages
    \item \texttt{yum install package} : installer
    \item \texttt{yum check-update} : mises \`a jour disponibles
    \item \texttt{yum update [package]} : mettre \`a jour
    \item \texttt{yum remove package} : d\'esinstaller
    \item \texttt{yum provides fichier} : chercher un paquet contenant un fichier
    \item \texttt{yum clean all} : vider le cache
    \item \texttt{yum makecache} : mettre \`a jour la liste des packages
    \item \texttt{yum grouplist} : lister les groupes de packages
    \item \texttt{yum groupinstall "nom\_groupe"} : installer un groupe
    \item \texttt{yum history} : historique des transactions
\end{itemize}

\subsection{Gestionnaire de paquets DNF}
Dandified Yum - successeur de YUM, plus performant.

Commandes similaires \`a YUM :
\begin{itemize}
    \item \texttt{dnf install/remove/update} : gestion des paquets
    \item \texttt{dnf search/info/provides} : recherche et informations
    \item \texttt{dnf list [available/installed/updates]} : listes
    \item \texttt{dnf groupinstall/groupremove} : gestion des groupes
    \item \texttt{dnf history} : historique
\end{itemize}

\subsection{D\'ep\^ots de paquetages}
\begin{itemize}
    \item Fichiers .repo dans /etc/yum.repos.d/
    \item \texttt{yum repolist all} : lister tous les d\'ep\^ots
    \item \texttt{yum repolist} : lister les d\'ep\^ots actifs
    \item \texttt{yum-config-manager --enable nom\_du\_d\'ep\^ot} : activer un d\'ep\^ot
\end{itemize}

\section{Gestion des Utilisateurs}

\subsection{Types de comptes utilisateurs}
\begin{itemize}
    \item \textbf{Super-utilisateur (root)} : UID = 0, tous les privil\`eges
    \item \textbf{Comptes syst\`emes} : UID $<$ 1000, services du syst\`eme
    \item \textbf{Comptes ordinaires} : UID $\geq$ 1000, utilisateurs normaux
\end{itemize}

\subsection{Fichiers d'enregistrement}
\begin{itemize}
    \item \textbf{/etc/passwd} : liste des utilisateurs avec descriptions
    \item \textbf{/etc/shadow} : mots de passe crypt\'es + dur\'ees de validit\'e
    \item \textbf{/etc/group} : d\'efinitions des groupes
    \item \textbf{/etc/gshadow} : mots de passe des groupes (rarement utilis\'e)
\end{itemize}

\subsection{Structure de /etc/passwd}
Format : \texttt{login:passwd:uid:gid:comment:home:shell}
\begin{itemize}
    \item login : nom de connexion
    \item passwd : x (indique pr\'esence de /etc/shadow)
    \item uid : num\'ero unique de l'utilisateur
    \item gid : num\'ero de groupe primaire
    \item comment : nom complet
    \item home : r\'epertoire personnel
    \item shell : interpr\'eteur de commandes
\end{itemize}

\subsection{Structure de /etc/shadow}
Format : \texttt{login:passwd:last:may:must:warn:inactif:disable}
\begin{itemize}
    \item passwd : mot de passe encod\'e
    \item last : jours depuis modification (depuis 1/1/1970)
    \item may : jours minimum avant changement
    \item must : jours maximum avant expiration
    \item warn : jours d'avertissement avant expiration
    \item inactif : p\'eriode d'inactivit\'e maximale
    \item disable : date de fermeture du compte
\end{itemize}

\subsection{Commandes de gestion des utilisateurs}

\subsubsection{Cr\'eation}
\begin{itemize}
    \item \texttt{useradd [options] nom\_utilisateur}
    \begin{itemize}
        \item \texttt{-u} : sp\'ecifier UID
        \item \texttt{-c} : commentaire (nom complet)
        \item \texttt{-d} : r\'epertoire personnel
        \item \texttt{-e} : date d'expiration (AAAA-MM-JJ)
        \item \texttt{-g} : groupe primaire
        \item \texttt{-G} : groupes secondaires
        \item \texttt{-m} : cr\'eer le r\'epertoire personnel
        \item \texttt{-s} : shell
        \item \texttt{-r} : cr\'eer un compte syst\`eme
    \end{itemize}
    \item \texttt{passwd utilisateur} : d\'efinir le mot de passe
\end{itemize}

\subsubsection{Modification}
\begin{itemize}
    \item \texttt{usermod [options] utilisateur}
    \begin{itemize}
        \item \texttt{-s} : changer le shell
        \item \texttt{-c} : modifier le commentaire
        \item \texttt{-aG} : ajouter \`a des groupes (sans \'ecraser)
        \item \texttt{-d -m} : d\'eplacer le home
        \item \texttt{-L} : verrouiller le compte
        \item \texttt{-U} : d\'everrouiller
        \item \texttt{-e} : expirer le compte
    \end{itemize}
\end{itemize}

\subsubsection{Suppression}
\begin{itemize}
    \item \texttt{userdel utilisateur} : supprimer (garde le home)
    \item \texttt{userdel -r utilisateur} : supprimer avec le home
    \item \texttt{userdel -rf utilisateur} : forcer m\^eme si connect\'e
\end{itemize}

\subsubsection{Gestion des mots de passe}
\begin{itemize}
    \item \texttt{passwd [options] utilisateur}
    \begin{itemize}
        \item \texttt{-d} : supprimer le mot de passe
        \item \texttt{-l} : verrouiller
        \item \texttt{-u} : d\'everrouiller
        \item \texttt{-e} : faire expirer
        \item \texttt{-n jours} : jours minimum avant changement
        \item \texttt{-x jours} : jours maximum avant expiration
        \item \texttt{-w jours} : d\'elai d'avertissement
        \item \texttt{-i jours} : d\'elai avant d\'esactivation
    \end{itemize}
    \item \texttt{chage [options] utilisateur} : modifier validit\'e du mot de passe
\end{itemize}

\subsection{Gestion des groupes}

\subsubsection{Cr\'eation et modification}
\begin{itemize}
    \item \texttt{groupadd [options] nom\_groupe}
    \begin{itemize}
        \item \texttt{-g GID} : sp\'ecifier GID
        \item \texttt{-r} : cr\'eer un groupe syst\`eme
    \end{itemize}
    \item \texttt{groupmod [options] nom\_groupe}
    \begin{itemize}
        \item \texttt{-g GID} : changer GID
        \item \texttt{-n nouveau\_nom} : renommer
    \end{itemize}
    \item \texttt{groupdel nom\_groupe} : supprimer
\end{itemize}

\subsubsection{Gestion des membres}
\begin{itemize}
    \item \texttt{gpasswd -a utilisateur groupe} : ajouter un membre
    \item \texttt{gpasswd -d utilisateur groupe} : retirer un membre
    \item \texttt{gpasswd -M utilisateur1,utilisateur2 groupe} : d\'efinir liste de membres
\end{itemize}

\subsection{Droits d'acc\`es}

\subsubsection{Permissions traditionnelles}
Trois cat\'egories : propri\'etaire (u), groupe (g), autres (o)
\begin{itemize}
    \item \textbf{r} (read) : lecture
    \item \textbf{w} (write) : \'ecriture
    \item \textbf{x} (execute) : ex\'ecution
\end{itemize}

\subsubsection{Modification des permissions}
\begin{itemize}
    \item \texttt{chmod [mode] fichier}
    \begin{itemize}
        \item Mode symbolique : \texttt{u+x}, \texttt{g-rw}, \texttt{o=r}, \texttt{a+r}
        \item Mode absolu : \texttt{755}, \texttt{644}, \texttt{600}, \texttt{777}
    \end{itemize}
    \item \texttt{chmod -R mode r\'epertoire} : r\'ecursif
\end{itemize}

\subsubsection{Changement de propri\'etaire}
\begin{itemize}
    \item \texttt{chown utilisateur[:groupe] fichier} : changer propri\'etaire
    \item \texttt{chgrp groupe fichier} : changer groupe
    \item \texttt{chown -R utilisateur:groupe r\'epertoire} : r\'ecursif
\end{itemize}

\subsection{Gestion des privil\`eges : sudo}
\begin{itemize}
    \item Fichier de configuration : /etc/sudoers
    \item \texttt{visudo} : \'editer de fa\c{c}on s\'ecuris\'ee
    \item Format : \texttt{utilisateur h\^ote=(utilisateur\_cible) commandes}
    \item Exemples :
    \begin{itemize}
        \item \texttt{jean ALL=(ALL:ALL) ALL} : tout faire
        \item \texttt{\%wheel ALL=(ALL) ALL} : groupe wheel
        \item \texttt{marie ALL=(root) NOPASSWD:/usr/sbin/reboot} : sans mot de passe
    \end{itemize}
\end{itemize}

\section{Gestion des Disques et Syst\`emes de Fichiers}

\subsection{Types de disques}
\begin{itemize}
    \item \textbf{HDD} : Disque dur m\'ecanique, \'economique, lent
    \item \textbf{SSD} : Disque \'electronique, rapide, silencieux
    \item \textbf{NVMe} : Tr\`es haute performance
\end{itemize}

\subsection{Partitionnement}

\subsubsection{MBR vs GPT}
\begin{itemize}
    \item \textbf{MBR} : Master Boot Record
    \begin{itemize}
        \item Limite : 4 partitions primaires maximum
        \item Taille max : 2.2 To par partition
        \item Compatible BIOS
    \end{itemize}
    \item \textbf{GPT} : GUID Partition Table
    \begin{itemize}
        \item Jusqu'\`a 128 partitions
        \item Taille tr\`es grande (jusqu'\`a 9.4 Zo)
        \item Compatible UEFI et BIOS
    \end{itemize}
\end{itemize}

\subsubsection{Organisation des disques sous Linux}
\begin{itemize}
    \item Format : /dev/sdXY ou /dev/nvme0n1pX
    \item Types de bus : hd (IDE), sd (SATA/SCSI/SSD), nvme (NVMe)
    \item UUID : identifiant unique universel
\end{itemize}

\subsubsection{Commandes d'information}
\begin{itemize}
    \item \texttt{lsblk} : afficher informations sur disques et partitions
    \item \texttt{blkid} : obtenir UUID
    \item \texttt{df -h} : taux de remplissage
    \item \texttt{du -h [chemin]} : espace occup\'e par une arborescence
\end{itemize}

\subsubsection{Cr\'eation de partitions}
\begin{itemize}
    \item \texttt{fdisk /dev/sdX} : manipuler table de partitions (interactif)
    \item \texttt{gdisk /dev/sdX} : pour GPT
    \item \texttt{parted} : mode interactif, script ou mixte
    \item Commandes fdisk : n (nouvelle), d (supprimer), t (type), w (sauvegarder), g (GPT)
\end{itemize}

\subsection{Syst\`emes de fichiers}

\subsubsection{D\'efinition}
Un syst\`eme de fichiers d\'efinit l'organisation d'un volume physique ou logique, permettant de stocker et organiser les informations dans des fichiers.

\subsubsection{Syst\`emes de fichiers GNU/Linux}
\begin{itemize}
    \item \textbf{ext2} : stable et robuste
    \item \textbf{ext3} : + journalisation
    \item \textbf{ext4} : volumes $>$ 1To, timestamps nanosec
    \item \textbf{XFS} : performance fichiers volumineux
    \item \textbf{Btrfs} : snapshots et compression
    \item \textbf{ZFS} : int\'egration stockage avanc\'ee
\end{itemize}

\subsubsection{Initialisation (formatage)}
\begin{itemize}
    \item \texttt{mkfs -t type partition}
    \item \texttt{mkfs.ext4 /dev/sda1}
    \item \texttt{lsblk -f} : v\'erifier le type
\end{itemize}

\subsection{Repr\'esentation des fichiers sur disque}

\subsubsection{Concepts}
\begin{itemize}
    \item \textbf{Bloc} : unit\'e de transfert entre disque et m\'emoire (souvent 4096 octets)
    \item \textbf{Inode} : descripteur d'un fichier contenant les m\'etadonn\'ees
\end{itemize}

\subsubsection{Inode}
Contient :
\begin{itemize}
    \item Type de l'inode (fichier, r\'epertoire)
    \item Propri\'etaire (UID, GID)
    \item Droits d'acc\`es
    \item Dates (atime, mtime, ctime)
    \item Taille
    \item Adresse : pointeurs sur les blocs du disque
\end{itemize}
\textbf{Important} : Les inodes ne contiennent pas les noms des fichiers.

\subsubsection{Liens}
\begin{itemize}
    \item \textbf{Lien physique} : \texttt{ln fichier lien}
    \begin{itemize}
        \item Partage le m\^eme inode
        \item Reste valide si fichier d\'eplac\'e
        \item Doit rester dans la m\^eme partition
    \end{itemize}
    \item \textbf{Lien symbolique} : \texttt{ln -s fichier lien}
    \begin{itemize}
        \item Pointe vers le nom du fichier
        \item Peut traverser les partitions
        \item Ressemble aux raccourcis
    \end{itemize}
\end{itemize}

\subsection{Montage d'une partition}

\subsubsection{Commandes}
\begin{itemize}
    \item \texttt{mount partition point\_de\_montage} : monter
    \item \texttt{mount -t type partition point} : sp\'ecifier type
    \item \texttt{mount -a} : monter tous les SF de /etc/fstab
    \item \texttt{umount point\_de\_montage} : d\'emonter
\end{itemize}

\subsubsection{Fichier /etc/fstab}
Configuration statique des montages. Structure :
\begin{verbatim}
périphérique  point_montage  type  options  dump  fsck
\end{verbatim}
\begin{itemize}
    \item dump : activer sauvegarde (1 ou 0)
    \item fsck : ordre de v\'erification (0, 1, 2, 3)
\end{itemize}

\subsection{Gestion de l'espace Swap}

\subsubsection{D\'efinition}
Zone de swap utilis\'ee lorsque la m\'emoire physique (RAM) est remplie. Les pages m\'emoire inactives sont d\'eplac\'ees dans la zone de swap.

\subsubsection{Configuration}
\begin{enumerate}
    \item Cr\'eer un fichier swap : \texttt{fallocate -l 1G /swapfile} ou \texttt{dd if=/dev/zero of=/swapfile bs=1M count=1024}
    \item Verrouiller permissions : \texttt{chmod 600 /swapfile}
    \item Marquer comme swap : \texttt{mkswap /swapfile}
    \item Activer : \texttt{swapon /swapfile}
    \item Rendre permanent : ajouter dans /etc/fstab
\end{enumerate}

\subsubsection{Commandes}
\begin{itemize}
    \item \texttt{swapon --show} : afficher swaps actifs
    \item \texttt{free -h} : \'etat m\'emoire et swap
    \item \texttt{swapoff /swapfile} : d\'esactiver
\end{itemize}

\section{S\'ecurit\'e sous Linux}

\subsection{ACLs : Liste de Contr\^ole d'Acc\`es}

\subsubsection{D\'efinition}
Les ACLs permettent d'\'etendre le nombre d'utilisateurs et de groupes ayant des droits sur un m\^eme fichier, au-del\`a des permissions traditionnelles.

\subsubsection{Commandes}
\begin{itemize}
    \item \texttt{getfacl fichier} : consulter les ACLs
    \item \texttt{setfacl -m u:utilisateur:rw fichier} : modifier ACLs
    \item \texttt{setfacl -m g:groupe:r fichier} : droits pour un groupe
    \item \texttt{setfacl -m o::- fichier} : refuser aux autres
    \item \texttt{setfacl -R -m u:utilisateur:rw r\'epertoire} : r\'ecursif
    \item \texttt{setfacl -b fichier} : supprimer toutes les ACLs
    \item \texttt{setfacl -x u:utilisateur fichier} : retirer permissions
\end{itemize}

\subsubsection{ACLs par d\'efaut}
\begin{itemize}
    \item \texttt{setfacl -m d:o::- r\'epertoire} : ACLs par d\'efaut
    \item S'appliquent aux nouveaux fichiers/r\'epertoires cr\'e\'es
    \item \texttt{setfacl -k r\'epertoire} : annuler ACLs par d\'efaut
\end{itemize}

\subsubsection{Masque}
Le masque limite les autorisations effectives pour les groupes et entr\'ees sp\'ecifiques.
\begin{itemize}
    \item \texttt{setfacl -m m::rwx fichier} : d\'efinir un masque
\end{itemize}

\subsection{Droits sp\'eciaux}

\subsubsection{Setuid}
Permet \`a un utilisateur non-root d'ex\'ecuter un programme avec les droits du propri\'etaire du programme.
\begin{itemize}
    \item \texttt{chmod u+s programme} : activer
    \item \texttt{chmod 4xxx programme} : mode absolu
    \item Affichage : 's' \`a la place de 'x' pour user
    \item 'S' si pas de droit x
\end{itemize}

\subsubsection{Setgid}
\begin{itemize}
    \item \textbf{Sur fichiers} : s'ex\'ecute avec les privil\`eges du groupe propri\'etaire
    \item \textbf{Sur r\'epertoires} : fichiers cr\'e\'es h\'eritent du groupe du r\'epertoire
    \item \texttt{chmod g+s r\'epertoire} : activer
    \item \texttt{chmod 2xxx r\'epertoire} : mode absolu
    \item Affichage : 's' \`a la place de 'x' pour groupe
\end{itemize}

\subsubsection{Sticky Bit}
Sur les r\'epertoires, emp\^eche les utilisateurs non-root de supprimer ou modifier les fichiers dont ils ne sont pas propri\'etaires.
\begin{itemize}
    \item \texttt{chmod +t r\'epertoire} : activer
    \item \texttt{chmod 1xxx r\'epertoire} : mode absolu
    \item Affichage : 't' \`a la place de 'x' pour other
    \item 'T' si pas de droit x
\end{itemize}

\subsection{SELinux : Security-Enhanced Linux}

\subsubsection{Introduction}
SELinux est une architecture de s\'ecurit\'e pour syst\`emes Linux permettant aux administrateurs de mieux contr\^oler les acc\`es au syst\`eme. C'est une couche suppl\'ementaire de s\'ecurit\'e.

\subsubsection{DAC vs MAC}
\begin{itemize}
    \item \textbf{DAC} (Discretionary Access Control) : contr\^ole d'acc\`es discr\'etionnaire
    \begin{itemize}
        \item Permissions classiques (r, w, x)
        \item Droits sp\'eciaux (setuid, setgid, sticky bit)
        \item ACLs
    \end{itemize}
    \item \textbf{MAC} (Mandatory Access Control) : contr\^ole d'acc\`es obligatoire
    \begin{itemize}
        \item Politique de s\'ecurit\'e \`a laquelle m\^eme root doit se soumettre
        \item Impl\'ementation : SELinux avec mod\`ele TE (Type Enforcement)
    \end{itemize}
\end{itemize}

\subsubsection{Concepts SELinux}
\begin{itemize}
    \item \textbf{Sujet} : entit\'e effectuant une action (processus, service, utilisateur)
    \item \textbf{Objet} : ressource sur laquelle une action est r\'ealis\'ee (fichier, r\'epertoire, socket, port)
    \item \textbf{R\`egles} : ensembles de r\`egles contr\^olant quelles actions un sujet peut effectuer sur un objet
    \item \textbf{Labeling} : chaque ressource poss\`ede une \'etiquette (contexte SELinux)
\end{itemize}

\subsubsection{Contexte SELinux}
Format : \texttt{user:role:type:level}
\begin{itemize}
    \item \textbf{user} : utilisateur SELinux (system\_u, unconfined\_u)
    \item \textbf{role} : r\^ole (object\_r pour objets)
    \item \textbf{type} : le plus important, d\'etermine les permissions
    \begin{itemize}
        \item httpd\_t : processus serveur web
        \item httpd\_sys\_content\_t : fichiers accessibles par Apache
        \item shadow\_t : fichier des mots de passe
    \end{itemize}
    \item \textbf{level} : niveau de sensibilit\'e (s0)
\end{itemize}

\subsubsection{Modes de fonctionnement}
\begin{itemize}
    \item \textbf{Enforcing} : mode strict, acc\`es restreints, refus des acc\`es non autoris\'es
    \item \textbf{Permissive} : mode d\'ebogage, r\`egles v\'erifi\'ees, erreurs enregistr\'ees mais acc\`es non bloqu\'e
    \item \textbf{Disabled} : d\'esactiv\'e, aucune restriction
\end{itemize}

\subsubsection{Commandes}
\begin{itemize}
    \item \texttt{getenforce} : mode actuel
    \item \texttt{setenforce enforcing} : passer en enforcing
    \item \texttt{setenforce permissive} : passer en permissif
    \item \texttt{sestatus} : \'etat g\'en\'eral
    \item \texttt{ls -Z fichier} : contexte d'un fichier
    \item \texttt{ps -Z} : contexte des processus
    \item \texttt{chcon -t type fichier} : changer contexte
    \item \texttt{restorecon fichier} : restaurer contexte par d\'efaut
    \item \texttt{semanage fcontext -l} : lister contextes par d\'efaut
\end{itemize}

\subsubsection{Configuration}
Fichier /etc/selinux/config :
\begin{verbatim}
SELINUX=enforcing|permissive|disabled
SELINUXTYPE=targeted
\end{verbatim}

\subsubsection{Exemples de types SELinux}
\begin{itemize}
    \item \texttt{httpd\_sys\_content\_t} : contenu accessible par Apache
    \item \texttt{device\_t} : fichiers associ\'es aux p\'eriph\'eriques
    \item \texttt{home\_dir\_t} : r\'epertoires personnels
    \item \texttt{var\_log\_t} : fichiers de journaux
    \item \texttt{shadow\_t} : fichiers de mots de passe
\end{itemize}

\section*{Conclusion}

Ce r\'esum\'e couvre les aspects essentiels de l'administration Unix/Linux, en excluant les parties sp\'ecifi\'ees (int\'egration r\'eseaux, LVM, RAID, sauvegarde/restauration, quotas). Les concepts pr\'esent\'es sont fondamentaux pour la gestion d'un syst\`eme Linux en environnement professionnel.

\end{document}

